%!TEX options=-shell-escape
%!TEX program=pdflatex

% Устанавливаем шрифт для всего документа
% Это нужно делать обязательно до загрузки memoir!
% \newcommand*{\memfontpack}{XCharter}
\newcommand*{\memfontpack}{tempora}
%!TEX root = graduate-work.tex
% Добавьте ссылку на файлы с текстом работы
% Можно использовать команды:
%   \input или \include
% Пример:
%    \input{mainfiles/1-section} или \include{mainfiles/2-section}
% Команда \input позволяет включить текст файла без дополнительной обработки
% Команда \include при включении файла добавляет до него и после него команду
% перехода на новую страницу. Кроме того, она позволяет компилировать каждый файл
% в отдельности, что ускоряет сборку проекта.
% ВАЖНО: команда \include не поддерживает включение файлов, в которых уже содержится команда \include,
% т.е. не возможен рекурсивный вызов \include
\newcommand*{\Source}{
    %!TEX root = ../graduate-work.tex
\phantomsection
\section*{Введение} 
\addcontentsline{toc}{section}{Введение}

В этой части надо описать предметную область, задачу из которой вы будете решать, объяснить её актуальность (почему надо что-то делать сейчас?).
Здесь же стоит ввести определения понятий, которые вам понадобятся в постановке задачи.

Данная система основывается на $\mathbf{\mathbb {NP}}$-трудной проблеме ... Однако, у неё имеется серьёзный недостаток ...

В этой работе рассматривается ...  предложенное в 1994 коду~\cite{Ivanov}. В этой работе ...

В настоящей работе исследуются вопросы ...

В \S~1 даётся определение ...


Для некоторых классов приводятся нижние оценки ... (теорема~\ref{theorem1}).
    %!TEX root = ../graduate-work.tex

\section{Постановка задачи}

Здесь надо максимально формально описать суть задачи, которую потребуется решить, так, чтобы можно было потом понять, в какой степени полученное в результате работы решение ей соответствует. Текст главы должен быть написан в стиле технического задания, т.е. содержать как описание задачи, так и некоторый набор требований к решению

\textbf{Алгоритм Решения Задачи.}
\begin{itemize}
\item[1.] Описать Задачу
\item[2.] Декомпозировать задачу $a=mH\in\mathcal
C,\;wt(e')=\lfloor\frac{d-1}2\rfloor;$
\item[3.] Решить каждую подзадачу и объединить решение.
\end{itemize}
    %!TEX root = ../graduate-work.tex

\section{Обзор существующих решений}

Здесь надо рассмотреть все существующие решения поставленной задачи, но не просто пересказать, в чем там дело, а оценить степень их соответствия тем ограничениям, которые были сформулированы в постановке задачи.

Информацию о группах автоморфизмов других известных кодов можно найти в~\cite{FirstArticle}.
    %!TEX root = ../graduate-work.tex

\section{Исследование и построение решения задачи}
Здесь надо декомпозировать большую задачу из постановки на подзадачи и продолжать этот процесс, пока подзадачи не станут достаточно простыми, чтобы их можно было бы решить напрямую (например, поставив какой-то эксперимент или доказав теорему) или найти готовое решение.

\begin{theorem}\label{theorem1}

\begin{itemize}
    \item[1)] Гомоморфный образ группы
    \item[2)] Что во имя коммунизма
    \item[3)] Изоморфен фактор-группе
    \item[4)] По ядру гомоморфизма
\end{itemize}
\end{theorem}
\begin{proof}
    Доказательство оставляем читателю.
\end{proof}
    %!TEX root = ../graduate-work.tex

\section[Опи\-са\-ние практи\-чес\-кой час\-ти с переносами]{Описание практической части с переносами}

Если в рамках работы писался какой-то код, здесь должно быть его описание: выбранный язык и библиотеки и мотивы выбора, архитектура, схема функционирования, теоретическая сложность алгоритма, характеристики функционирования (скорость/память).
    %!TEX root = ../graduate-work.tex

\section{Заключение}

Здесь надо перечислить все результаты, полученные в ходе работы. Из текста должно быть понятно, в какой мере решена поставленная задача.

В каком-то смысле в заключении повторяются введе
}


% Информация о годе выполнения работы
\def\Year{%
    % 2006%
    2026%     % Текущий год
}

% Укажите тип работы
% Например:
%     Выпускная квалификационная работа,
%     Магистерская диссертация,
%     Курсовая работа, реферат и т.п.
\def\WorkType{%
    Численное решение модели реального делового цикла для закрытой экономики%
    % Магистерская диссертация%
    % Курсовая работа%
    % Реферат%
    %Дипломная работа%
}

% Название работы
%%%%%%%%%%% ВНИМАНИЕ! %%%%%%%%%%%%%%%%
% В МГУ ОНО ДОЛЖНО В ТОЧНОСТИ
% СООТВЕТСТВОВАТЬ ВЫПИСКЕ ИЗ ПРИКАЗА
% УТОЧНИТЕ НАЗВАНИЕ В УЧЕБНОЙ ЧАСТИ
\def\Title{%
    Численное решение модели реального делового цикла для закрытой экономики%
}


% Имя автора работы
\def\Author{%
    Козлов Кирилл Андреевич%
}

% Информация о научном руководителе
%% Фамилия Имя Отчество%
\def\SciAdvisor{%
    Богомолов Сергей Владимирович%
}
%% В формате: И.~О.~Фамилия%
\def\SciAdvisorShort{%
    С.~В.~Богомолов%
}
%% должность научного руководителя
\def\Position{%
    профессор%
    % доцент%
    % старший преподаватель%
    % преподаватель%
    % ассистент%
    % ведущий научный сотрудник%
    % старший научный сотрудник%
    % научный сотрудник%
    % младший научный сотрудник%
}
%% учёная степень научного руководителя
\def\AcademicDegree{%
    д.ф.-м.н.%
    % д.т.н.%
    % к.ф.-м.н.%
    % к.т.н.%
    % без степени%
}

% Информация об организации, в которой выполнена работа
%% Город
\def\Place{%
    Москва%
}
%% Университет
\def\Univer{%
    Московский государственный университет имени М.~В.~Ломоносова%
}
%% Факультет
\def\Faculty{%
    Факультет вычислительной математики и кибернетики%
}
%% Кафедра    
\def\Department{%
    Кафедра вычислительных методов%
}     

%%%% Переключите статус документа для отладки
%%%% В режиме draft документ собирается очень быстро
%%%% и выводится полезная информация о том
%%%% какие строки вылезают за границы документа, что удобно для борьбы с ними
\def\Status{%
    draft%
    % final%
}

%%%% Включает и выключает подпись <<С текстом работы ознакомлен>>
\def\EnableSign{%
    % true%
}

\documentclass[
    a4paper,
    article,
    oneside,
    onecolumn,
    12pt,
    extrafontsizes,
    openright
]{memoir}

%Полезный условный оператор
\usepackage{ifthen}

%%%%%%% ВНИМАНИЕ! ИЗМЕНЕНИЕ ПОРЯДКА ПОДКЛЮЧЕНИЯ СТИЛЕЙ МОЖЕТ НАРУШИТЬ РАБОТУ ШАБЛОНА %%%%%%%
%Настройка шрифтов
\input{styles/fonts}
%Обеспечиваем работу с математикой
\input{styles/math}
%Настраиваем дизайн текста (отступы, интервалы и пр.)
%!TEX root = ../diploma.tex

%Отступ в первом абзаце раздела
\usepackage{indentfirst}
%Точка после номера раздела
\usepackage{misccorr}
% Настройка списков с пользовательскими метками (label=...)
\usepackage{enumitem}

\usepackage{hyperref}
%Тонкое оформление перекрёстных ссылок
\usepackage[russian]{cleveref}
\usepackage[dvipsnames, table, hyperref]{xcolor}
\usepackage{csquotes}
\usepackage{hyphenat}
% \hyphenation{ма-те-ма-ти-ка вос-ста-нав-ли-вать}


% Принудительное уменьшение размера заголовков глав
\usepackage{titlesec}

% Уменьшение размера заголовков секций
\titleformat{\chapter}[hang]{\large\bfseries}{\thechapter}{1em}{}
\titleformat{\section}{\large\bfseries}{}{0pt}{}
\titleformat{\subsection}{\normalsize\bfseries}{}{0pt}{}
\titleformat{\subsubsection}{\normalsize\itshape}{}{0pt}{}
% Для ненумерованных глав (звёздочка)
\titleformat{name=\chapter,numberless}[hang]{\large\bfseries}{}{0pt}{} % Если вы хотите ещё меньший шрифт, замените \large на \normalsize или \small.
% Настройка отступов до и после заголовка
\titlespacing*{\chapter}{0pt}{0pt}{20pt}

% Скрываем номера section/subsection/subsubsection в оглавлении
\makeatletter
\let\origl@section\l@section
\renewcommand*{\l@section}[2]{%
  \begingroup
  \let\numberline\@gobble
  \origl@section{#1}{#2}%
  \endgroup
}
\let\origl@subsection\l@subsection
\renewcommand*{\l@subsection}[2]{%
  \begingroup
  \let\numberline\@gobble
  \origl@subsection{#1}{#2}%
  \endgroup
}
\let\origl@subsubsection\l@subsubsection
\renewcommand*{\l@subsubsection}[2]{%
  \begingroup
  \let\numberline\@gobble
  \origl@subsubsection{#1}{#2}%
  \endgroup
}
\makeatother

%Настройка дизайна страницы
\input{styles/page}
%Обеспечиваем работу с медиа-контентом (с рисунками)
\input{styles/media}
%Обеспечиваем работу с таблицами
\input{styles/table}
%Обеспечиваем оформление алгоритмов и листингов
\input{styles/algorithms}

% Оформление библиографии (BibTeX, нумерованные ссылки)
\bibliographystyle{unsrt}

% Команды шаблона
\input{styles/commands}

% Пользовательские команды
\input{commands}


\begin{document}
    %Титульная страница
    \input{styles/title}

    %Вставляем оглавление
    \clearpage
    \tableofcontents*
    \clearpage

    %Вставляем основное содержание
    \phantomsection
\section*{Введение}
\addcontentsline{toc}{chapter}{Введение}

Модели общего динамического равновесия занимают важное место в современной макроэкономической теории \cite{heermaussner2009dynamic}. Они позволяют исследовать совместную динамику выпуска, потребления, капитала и труда, а также анализировать механизм межвременного выбора экономических агентов. Одной из наиболее известных конструкций такого типа является модель реального делового цикла (Real Business Cycle, RBC), в которой колебания основных макроэкономических переменных объясняются реакцией экономики на реальные шоки при рациональном поведении домохозяйств и фирм \cite{kydland1982time}. Дальнейшим развитием этого подхода стали неокейнсианские DSGE-модели (Dynamic Stochastic General Equilibrium), которые дополнительно учитывают номинальные жёсткости (цен и заработных плат), несовершенную конкуренцию и денежно-кредитную политику \cite{smets2007shocks, christiano2005nominal}.

Модели реального делового цикла (RBC) и их неокейнсианские расширения (DSGE) являются стандартным инструментом макроэкономического анализа, используемым центральными банками, инвестиционными фондами и министерствами экономики для прогнозирования, стресс-тестирования и оценки последствий политических решений \cite{motto2025chapter, ristiniemi2025chapter}. Согласно опросу 22 центральных банков, DSGE-модели служат основной или вспомогательной моделью для большинства из них \cite{zamrazilova2025role}. Однако даже детерминированная RBC-модель приводит к нелинейной системе разностных уравнений, аналитическое решение которой невозможно. Традиционные локальные методы (линеаризация) дают приемлемую точность лишь в малой окрестности стационарного состояния \cite{judd1998numerical}, тогда как реальные шоки могут быть значительными. Поэтому актуальной задачей является реализация и сопоставление глобальных численных методов (например, метода стрельбы) с локальными подходами для построения переходных траекторий при больших отклонениях от равновесия \cite{lipton1983multiple}.

Целью выпускной квалификационной работы является сопоставление глобального (метод стрельбы) и локального (линеаризация) подходов к решению детерминированной модели реального делового цикла для закрытой экономики \cite{judd1998numerical}, а также анализ её стационарного состояния и переходной динамики \cite{heermaussner2009dynamic}. Дополнительно модель расширяется на стохастический случай для анализа импульсных откликов на шоки совокупной факторной производительности и денежно-кредитной политики \cite{sims2002solving, blanchardkahn1980}.

Объектом исследования является детерминированная и стохастическая макроэкономическая модель закрытой экономики типа RBC. Предметом исследования выступают стационарное состояние модели, переходная динамика капитала, потребления, выпуска и труда, а также численные методы построения устойчивой траектории \cite{judd1998numerical, fernandez2015solution}.

Работа состоит из введения, трёх глав, заключения и списка литературы. В первой главе формально ставятся задачи домохозяйства и фирмы, выводится необходимое условие экстремума, исследуется стационарное состояние и формируется итоговая система уравнений детерминированной модели в динамической форме. Во второй главе рассматриваются существующие численные подходы к решению RBC-моделей: даётся классификация на локальные и глобальные методы, обосновывается выбор методов для данной работы, излагается алгоритм построения траектории при фиксированном начальном уровне потребления, метод Ньютона для нахождения управляющей переменной и приводится псевдокод метода стрельбы. В третьей главе представлены результаты численных экспериментов: фазовые траектории в плоскостях $(K, C)$ и $(K, L)$, трёхмерная траектория, а также сравнение глобального и локального решений. В четвёртой главе формулируются основные выводы и обсуждаются перспективы дальнейших исследований.
    \chapter{Описание модели}

В данной главе представлена стандартная модель реального делового цикла (Real Business Cycle, RBC) с дискретным временем. Экономика состоит из репрезентативного домохозяйства, репрезентативной фирмы и государства (которое для простоты отсутствует). Модель описывает поведение экономических агентов в условиях совершенной конкуренции, отсутствия внешних эффектов и полной предопределённости (детерминированная версия). Все переменные измеряются в реальном выражении.

\section{Экономическая среда}

Время дискретно, \(t = 0, 1, 2, \dots, \infty\). В каждом периоде домохозяйство принимает решения о потреблении \(C_t\) и предложении труда \(L_t\). Фирма производит однородный конечный продукт \(Y_t\) с использованием капитала \(K_t\) и труда \(L_t\) по технологии Кобба–Дугласа. Капитал накапливается в соответствии с законом движения, включающим амортизацию \(\delta \in (0,1)\). Технологический прогресс отсутствует, уровень совокупной факторной производительности (TFP) \(A\) постоянен.

\section{Задача домохозяйства}

Домохозяйство максимизирует дисконтированную сумму мгновенных полезностей, зависящих от потребления и труда:


\[
\max_{\{C_t, L_t\}_{t=0}^\infty} \sum_{t=0}^\infty \beta^t u(C_t, L_t),
\]

где \(\beta \in (0,1)\) – субъективный дисконт-фактор. Функция полезности задаётся в сепарабельном виде с постоянной эластичностью замещения (CRRA) по потреблению и линейным отрицательным вкладом труда:

\[
u(C_t, L_t) = \frac{C_t^{1-\sigma} - 1}{1-\sigma} - \varphi \frac{L_t^{1+\psi}}{1+\psi},
\]

где \(\sigma > 0\) – коэффициент относительного неприятия риска (обратная эластичность межвременного замещения), \(\varphi > 0\) – параметр, отражающий «ненависть к труду», \(\psi > 0\) – обратная эластичность предложения труда по Фришу.

Домохозяйство владеет начальным запасом капитала \(K_0\) и получает доход от сдачи капитала в аренду по ставке \(R_t\) и от продажи труда по ставке заработной платы \(w_t\). Доход используется на потребление и инвестиции \(I_t\). Бюджетное ограничение в каждом периоде имеет вид:

\[
C_t + I_t = w_t L_t + R_t K_t.
\]

Капитал изменяется согласно закону движения:

\[
K_{t+1} = I_t + (1-\delta)K_t.
\]

Домохозяйство выбирает \(\{C_t, L_t, K_{t+1}\}_{t=0}^\infty\), максимизируя полезность при заданных начальных условиях и бюджетных ограничениях.

\section{Задача фирмы}

Фирма производит продукт по технологии Кобба–Дугласа с постоянной отдачей от масштаба:

\[
Y_t = A K_t^\alpha L_t^{1-\alpha}, \quad \alpha \in (0,1),
\]

где \(A > 0\) – уровень TFP. Фирма максимизирует прибыль в каждом периоде:

\[
\max_{K_t, L_t} \; A K_t^\alpha L_t^{1-\alpha} - w_t L_t - R_t K_t.
\]

В силу совершенной конкуренции и отсутствия издержек на корректировку, условия первого порядка дают стандартные выражения для факторных цен:

\[
R_t = \alpha \frac{Y_t}{K_t}, \qquad w_t = (1-\alpha) \frac{Y_t}{L_t}.
\]

\section{Условия равновесия}

Равновесие в модели определяется как набор последовательностей \(\{C_t, L_t, K_{t+1}, Y_t, R_t, w_t\}_{t=0}^\infty\), удовлетворяющих:

\begin{itemize}
  \item задаче домохозяйства (условия оптимальности),
  \item задаче фирмы (предельные производительности факторов),
  \item бюджетному ограничению домохозяйства (которое с учётом закона движения капитала превращается в ресурсное ограничение экономики):
  \[
  K_{t+1} = Y_t - C_t + (1-\delta)K_t,
  \]
  \item начальному условию \(K_0\) (задано экзогенно)
\end{itemize}

\section{Необходимое условие экстремума}

Выпишем лагранжиан задачи домохозяйства:

\[
\mathcal{L} = \sum_{t=0}^\infty \beta^t \left[ \frac{C_t^{1-\sigma}-1}{1-\sigma} - \varphi \frac{L_t^{1+\psi}}{1+\psi} + \lambda_t \bigl( w_t L_t + R_t K_t - C_t - K_{t+1} + (1-\delta)K_t \bigr) \right],
\]

где \(\lambda_t\) – множитель Лагранжа для бюджетного ограничения в периоде \(t\).

Условия первого порядка (по \(C_t\), \(L_t\), \(K_{t+1}\)):

\begin{align}
\frac{\partial \mathcal{L}}{\partial C_t} &= 0 \quad \Rightarrow \quad \beta^t C_t^{-\sigma} - \lambda_t = 0, \tag{1}\\[4pt]
\frac{\partial \mathcal{L}}{\partial L_t} &= 0 \quad \Rightarrow \quad -\beta^t \varphi L_t^{\psi} + \lambda_t w_t = 0, \tag{2}\\[4pt]
\frac{\partial \mathcal{L}}{\partial K_{t+1}} &= 0 \quad \Rightarrow \quad -\lambda_t + \lambda_{t+1}\bigl(R_{t+1} + 1 - \delta\bigr) = 0. \tag{3}
\end{align}

Из (1) и (3) получаем уравнение Эйлера для потребления:

\[
C_t^{-\sigma} = \beta \, C_{t+1}^{-\sigma} \bigl( R_{t+1} + 1 - \delta \bigr).
\]

С учётом предельной производительности капитала \(R_{t+1} = \alpha Y_{t+1}/K_{t+1}\) уравнение принимает вид:

\[
\frac{C_{t+1}}{C_t} = \beta^{1/\sigma} \bigl( \alpha A K_{t+1}^{\alpha-1} L_{t+1}^{1-\alpha} + 1 - \delta \bigr)^{1/\sigma}.
\]

Из (1) и (2) выводится условие внутридневного выбора между потреблением и трудом (условие оптимального предложения труда):

\[
\varphi L_t^{\psi} = w_t \, C_t^{-\sigma}.
\]

Используя выражение для заработной платы \(w_t = (1-\alpha) A K_t^\alpha L_t^{-\alpha}\), получаем:

\[
\varphi L_t^{\psi} = (1-\alpha) A \left(\frac{K_t}{L_t}\right)^\alpha C_t^{-\sigma}.
\]

Таким образом, система условий первого порядка совместно с ресурсным ограничением и законом движения капитала полностью описывает динамику модели.

\section{Стационарное состояние (Steady State)}

  Стационарным состоянием (steady state) модели называется такая траектория, на которой все переменные остаются постоянными во времени. Иными словами, выполняются условия:
  \[
  K_{t+1}=K_t=K^*, \qquad 
  C_{t+1}=C_t=C^*, \qquad 
  L_{t+1}=L_t=L^*,
  \]
  и, как следствие, $Y_{t+1}=Y_t=Y^*$, $R_{t+1}=R_t=R^*$, $w_{t+1}=w_t=w^*$, $I_{t+1}=I_t=I^*$.

  Вывод аналитических выражений для стационарных значений осуществляется последовательно, начиная с уравнения Эйлера.


  Уравнение Эйлера для потребления в стационарном режиме принимает вид:
  \begin{equation}
  1 = \beta \bigl(1 + R^* - \delta\bigr).
  \end{equation}
  Отсюда немедленно получаем:
  \begin{equation}
  R^* = \frac{1}{\beta} - 1 + \delta. \tag{3.1}
  \end{equation}
  Таким образом, стационарная доходность капитала полностью определяется параметрами дисконтирования $\beta$ и выбытия $\delta$.


  Закон движения капитала в стационарном состоянии даёт:
  \[
  K^* = Y^* - C^* + (1-\delta)K^* \quad\Longrightarrow\quad Y^* = C^* + \delta K^*.
  \]

  С другой стороны, из определения доходности капитала $R^* = \alpha Y^*/K^*$ имеем $Y^* = \frac{R^*}{\alpha}K^*$. Приравнивая два выражения для $Y^*$, находим:
  \[
  \frac{R^*}{\alpha}K^* = C^* + \delta K^* \quad\Longrightarrow\quad 
  \frac{R^*}{\alpha} = \frac{C^*}{K^*} + \delta.
  \]

  Введём обозначение $v = C^*/K^*$ для отношения стационарного потребления к капиталу. Тогда:
  \begin{equation}
  v = \frac{C^*}{K^*} = \frac{R^*}{\alpha} - \delta. \tag{3.2}
  \end{equation}


  Для нахождения $L^*$ используем условие внутрипериодной оптимальности, связывающее заработную плату, потребление и труд:
  \[
  \phi C^* (L^*)^{\psi} = w^*,
  \]
  а также выражение для заработной платы через предельный продукт труда:
  \[
  w^* = (1-\alpha)\frac{Y^*}{L^*}.
  \]

  Подставляя $Y^* = \frac{R^*}{\alpha}K^*$ и $C^* = v K^*$, получаем:
  \[
  \phi v K^* (L^*)^{\psi} = (1-\alpha)\frac{R^*}{\alpha}\frac{K^*}{L^*}.
  \]

  Сокращая $K^*$ (положительный множитель) и умножая обе части на $L^*$, приходим к уравнению:
  \[
  \phi v (L^*)^{\psi+1} = \frac{1-\alpha}{\alpha}R^*.
  \]

  Отсюда выражается стационарный уровень труда:
  \begin{equation}
  L^* = \left( \frac{(1-\alpha)R^*}{\alpha \phi v} \right)^{\frac{1}{\psi+1}}. \tag{3.3}
  \end{equation}


  Теперь найдём $K^*$, используя производственную функцию Кобба–Дугласа:
  \[
  Y^* = A (K^*)^{\alpha} (L^*)^{1-\alpha}.
  \]

  С другой стороны, из $R^* = \alpha Y^*/K^*$ имеем $Y^* = \frac{R^*}{\alpha}K^*$. Приравнивая, получаем:
  \[
  \frac{R^*}{\alpha}K^* = A (K^*)^{\alpha} (L^*)^{1-\alpha}.
  \]

  Разрешая относительно $K^*$:
  \[
  (K^*)^{1-\alpha} = \frac{\alpha A}{R^*} (L^*)^{1-\alpha}
  \quad\Longrightarrow\quad
  K^* = \left( \frac{\alpha A}{R^*} \right)^{\frac{1}{1-\alpha}} L^*.
  \]

  Более удобная форма (явно выделяющая $L^*$):
  \begin{equation}
  K^* = \left( \frac{\alpha A (L^*)^{1-\alpha}}{R^*} \right)^{\frac{1}{1-\alpha}}. \tag{3.4}
  \end{equation}

  Наконец, стационарное потребление находится по определению $v$:
  \begin{equation}
  C^* = v K^*. \tag{3.5}
  \end{equation}


  Таким образом, стационарное состояние модели полностью описывается следующей системой аналитических выражений, вычисляемых в строгой последовательности:

  \[
  \begin{aligned}
  R^* &= \frac{1}{\beta} - 1 + \delta, \\[4pt]
  v &= \frac{R^*}{\alpha} - \delta, \\[4pt]
  L^* &= \left( \frac{(1-\alpha)R^*}{\alpha \phi v} \right)^{\frac{1}{\psi+1}}, \\[4pt]
  K^* &= \left( \frac{\alpha A (L^*)^{1-\alpha}}{R^*} \right)^{\frac{1}{1-\alpha}}, \\[4pt]
  C^* &= v K^*, \\[4pt]
  Y^* &= A (K^*)^{\alpha} (L^*)^{1-\alpha} = \frac{R^*}{\alpha} K^*, \\[4pt]
  I^* &= \delta K^*, \\[4pt]
  w^* &= (1-\alpha)\frac{Y^*}{L^*}.
  \end{aligned}
  \]

  Все величины выражаются исключительно через экзогенные параметры модели $\beta$, $\alpha$, $\delta$, $\phi$, $\psi$, $A$, что позволяет проводить численные расчёты и анализировать чувствительность равновесия к изменению параметров.
  
    \chapter{Метод стрельбы}
        
Для нахождения устойчивой переходной траектории используется метод стрельбы. Его идея состоит в следующем: при фиксированном начальном значении капитала $K_0$ подбирается такое начальное значение потребления $C_0$, при котором траектория системы на конечном горизонте максимально приближается к стационарному состоянию.

На каждом шаге алгоритма выполняются следующие действия:
\begin{enumerate}[label=\arabic*)]
    \item по заданным $K_t$ и $C_t$ вычисляется $L_t$;
    \item определяется выпуск $Y_t$ и инвестиции $I_t=Y_t-C_t$;
    \item по уравнению накопления вычисляется капитал $K_{t+1}$;
    \item численно решается нелинейное уравнение относительно $C_{t+1}$:
    \begin{equation}
    C_{t+1}=\beta\bigl(1+R(K_{t+1},C_{t+1})-\delta\bigr)C_t;
    \end{equation}
    \item процесс повторяется до достижения заданного горизонта моделирования.
\end{enumerate}

Подбор начального потребления осуществляется по минимуму функционала, измеряющего отклонение терминального состояния от стационарной точки.
    
\section{Построение траектории при фиксированном \texorpdfstring{$C_0$}{C0}}

    Пусть заданы начальные значения $K_0$ и $C_0$. Для каждого периода $t=0,1,\dots,T-1$ траектория строится итерационно. Сначала по внутрипериодному условию вычисляется труд:
    \[
    L_t=\left(\frac{(1-\alpha)A}{\phi}\frac{K_t^\alpha}{C_t}\right)^{\frac{1}{\psi+\alpha}}.
    \]
    Затем определяются выпуск и связанные с ним величины:
    \(Y_t=A K_t^\alpha L_t^{1-\alpha}\),
    \(R_t=\alpha\frac{Y_t}{K_t}\),
    \(I_t=Y_t-C_t\).
    После этого капитал следующего периода находится из закона накопления:
    \(K_{t+1}=Y_t-C_t+(1-\delta)K_t\).

    Значение $C_{t+1}$ определяется из уравнения Эйлера и находится как корень одномерного нелинейного уравнения
    \[
    f(c)=c-\beta\bigl(1+R(K_{t+1},c)-\delta\bigr)C_t=0,
    \]
    где \(c=C_{t+1}\).
    При этом функция доходности на шаге \(t+1\) задается через
    \(Y_{t+1}=A K_{t+1}^\alpha L_{t+1}^{1-\alpha}\),
    \(L_{t+1}=\left(\frac{(1-\alpha)A}{\phi}\frac{K_{t+1}^\alpha}{c}\right)^{\frac{1}{\psi+\alpha}}\),
    \(R_{t+1}=\alpha\frac{Y_{t+1}}{K_{t+1}}\).
    После нахождения \(C_{t+1}\) выполняется переход к периоду \(t+1\), и процедура повторяется до горизонта \(T\).
\section{Метод Ньютона для нахождения \texorpdfstring{$C_{t+1}$}{C(t+1)}}

    После вычисления \(K_{t+1}\) значение \(C_{t+1}\) определяется из неявного уравнения Эйлера:
    \[
    f(c)=c-\beta\bigl(1+R(K_{t+1},c)-\delta\bigr)C_t=0,\qquad c=C_{t+1}.
    \]
    Для фиксированного \(K=K_{t+1}\) используем обозначения
    \(L(K,c)=\left(\frac{(1-\alpha)A}{\phi}\frac{K^\alpha}{c}\right)^{\frac{1}{\psi+\alpha}}\),
    \(Y(K,c)=A K^\alpha L(K,c)^{1-\alpha}\),
    \(R(K,c)=\alpha\frac{Y(K,c)}{K}\).

    Производная \(f'(c)\) вычисляется аналитически. Из
    \(\frac{\partial L}{\partial c}=-\frac{1}{\psi+\alpha}\frac{L}{c}\),
    \(\frac{\partial Y}{\partial c}=-\frac{1-\alpha}{\psi+\alpha}\frac{Y}{c}\)
    получаем
    \(\frac{\partial R}{\partial c}=-\frac{1-\alpha}{\psi+\alpha}\frac{R}{c}\), а значит
    \[
    f'(c)=1-\beta C_t\frac{\partial R(K,c)}{\partial c}
    =1+\beta C_t\frac{1-\alpha}{\psi+\alpha}\frac{R(K,c)}{c}.
    \]

    Тогда итерации Ньютона записываются в виде
    \[
    c^{(n+1)}=
    c^{(n)}-\frac{f(c^{(n)})}{f'(c^{(n)})}
    =
    c^{(n)}-
    \frac{c^{(n)}-\beta\bigl(1+R(K_{t+1},c^{(n)})-\delta\bigr)C_t}
    {1+\beta C_t \frac{1-\alpha}{\psi+\alpha}\frac{R(K_{t+1},c^{(n)})}{c^{(n)}}}.
    \]
    После достижения критерия сходимости принимается \(C_{t+1}=c^{(n)}\).

\section{Псевдокод метода стрельбы}
    \begin{enumerate}[label=\arabic*)]
        \item Задать параметры модели $(\alpha,\beta,\delta,\phi,\psi,A)$.
        \item Вычислить стационарное состояние $(K^*,C^*,L^*)$.
        \item Зафиксировать $K_0$.
        \item Задать начальный интервал для потребления: $C_0\in[C_{\min},C_{\max}]$.
        \item Повторять внешнюю бисекцию по $C_0$, пока $|C_{\max}-C_{\min}|\ge\varepsilon_{\text{outer}}$.
        \item На каждой итерации положить $C_0=\frac{C_{\min}+C_{\max}}{2}$.
        \item Построить траекторию на горизонте $t=0,1,\dots,T-1$.
        \item Для каждого $t$ вычислить \(L_t=\left(\frac{(1-\alpha)A}{\phi}\frac{K_t^\alpha}{C_t}\right)^{\frac{1}{\psi+\alpha}}\).
        \item Для каждого $t$ вычислить \(Y_t=A K_t^\alpha L_t^{1-\alpha}\).
        \item Для каждого $t$ вычислить \(K_{t+1}=Y_t-C_t+(1-\delta)K_t\).
        \item Для каждого $t$ найти $C_{t+1}$ методом Ньютона из \(f(c)=c-\beta\bigl(1+R(K_{t+1},c)-\delta\bigr)C_t=0\).
        \item По завершении траектории вычислить терминальное отклонение $g(C_0)=K_T-K^*$.
        \item Если $g(C_0)>0$, положить $C_{\min}=C_0$.
        \item Если $g(C_0)\le0$, положить $C_{\max}=C_0$.
        \item После сходимости внешней бисекции принять найденное $C_0^*$.
        \item Построить итоговую saddle-path траекторию $(K_t,C_t,L_t)$.
    \end{enumerate}

    \chapter{Численное решение}

В данной главе представлены результаты вычислительных экспериментов для модели реального делового цикла, реализованных на основе построенного в предыдущей главе численного алгоритма. Основная цель главы --- показать, как построенное глобальное решение (метод стрельбы) описывает динамику модели и чем оно отличается от локального приближения.

В рамках практической части работы выполнены следующие шаги:
\begin{enumerate}
    \item реализована численная схема решения детерминированной RBC-модели с подбором начального потребления \(C_0\) методом стрельбы;
    \item для каждого шага по времени реализовано решение неявного уравнения Эйлера для \(C_{t+1}\) методом Ньютона;
    \item рассчитано стационарное состояние модели и построены переходные траектории из разных начальных условий;
    \item построены фазовые графики в координатах \((K_t, C_t)\) и \((K_t, L_t)\);
    \item построены трехмерные траектории для наглядного анализа выхода экономики к стационарной точке;
    \item выполнено визуальное сравнение локального и глобального решений для DSGE-постановки.
\end{enumerate}

Во всех экспериментах использовались параметры:
\[
\beta = 0.95,\quad
\alpha = 0.33,\quad
\delta = 0.1,\quad
\sigma = 1,\quad
\phi = 0.75,\quad
\psi = 0.35.
\]


\section{Фазовые траектории (K, C) и (K, L)}

На рис.~\ref{fig:rbc_phase_kc} и~\ref{fig:rbc_phase_kl} приведены фазовые траектории для детерминированной модели при различных начальных условиях А - уровня технического прогресса. Для экспериментов были выбраны значения А: $0.95, 1.0, 1.05$, вычислены соответствующие стационарные состояния и построены фазовыве траектории с помощью метода стрельбы слева и справа от стационарного состояния. Первый график показывает совместную динамику капитала и потребления \((K_t, C_t)\), второй --- капитала и труда \((K_t, L_t)\).

\begin{figure}[htbp]
    \centering
    \includegraphics[width=0.9\textwidth]{kc.png}
    \caption{Фазовая траектория детерминированной RBC-модели в координатах \((K_t, C_t)\)}
    \label{fig:rbc_phase_kc}
\end{figure}

\begin{figure}[htbp]
    \centering
    \includegraphics[width=0.9\textwidth]{kl.png}
    \caption{Фазовая траектория детерминированной RBC-модели в координатах \((K_t, L_t)\)}
    \label{fig:rbc_phase_kl}
\end{figure}

\FloatBarrier

Полученные графики подтверждают наличие устойчивой переходной траектории к стационарному состоянию при любых значениях параметра А. Подбор \(C_0\) методом стрельбы обеспечивает сходящуюся динамику без взрывного поведения переменных.

Экспериментально доказаны теорические основы макроэкономики: при росте капитала $K$ потребление $C$ у домохозяйства растёт, а труд $L$ падает. 

\section{Трёхмерная траектория}

Для дополнительной проверки устойчивости построены трехмерные седловые траектории из двух различных начальных условий по капиталу: $K_0 = 10*K^*$ и $K_0 = 0.4*K^*$ соответственно. На рис.~\ref{fig:rbc_3d_first} и~\ref{fig:rbc_3d_second} показаны соответствующие переходные траектории. 
Красными точками на этих двух графиках отмечено стационарное состояние(одинаковое для обоих графиков), зелёными - успешная стартовая точка, из которой метод стрельбы сошёлся к стационарному состоянию. Слегка прозрачными линиями - неудачные траектории, которые не сошлись к стационарному состоянию. 

\begin{figure}[htbp]
    \centering
    \includegraphics[width=0.9\textwidth]{rbc_3d.png}
    \caption{Трехмерная седловая траектория при $K_0 = 10*K^*$}
    \label{fig:rbc_3d_first}
\end{figure}

\begin{figure}[htbp]
    \centering
    \includegraphics[width=0.9\textwidth]{rbc_3d_2.png}
    \caption{Трехмерная седловая траектория при $K_0 = 0.4*K^*$}
    \label{fig:rbc_3d_second}
\end{figure}

\FloatBarrier

Обе траектории сходятся к одной и той же стационарной точке, что является ожидаемым результатом для седлово-устойчивой динамической системы в окрестности равновесия. На графиках виден разброс, которые может возникнуть при неудачном выборе начальное потребления $C_0$

\section{Сравнение локального и глобального решения DSGE}

В DSGE-модели решением является распределение траекторий макропеременных, а не одна детерминированная траектория. На рис.~\ref{fig:local_global_dsge} и~\ref{fig:local_global_dsge_2} приведено сравнение локального и глобального подходов по всем десяти эндогенным переменным модели. Для удобства восприятия они разбиты на два рисунка: реальные величины и труд (рис.~\ref{fig:local_global_dsge}) и цены и факторы (рис.~\ref{fig:local_global_dsge_2}). По оси Y указаны проценты отклонения от стационарного состояния (за исключением прибыли \(\mathit{Pr}\), для которой используется отклонение в уровнях), по оси X --- периоды симуляции.

\begin{figure}[htbp]
    \centering
    \includegraphics[width=\textwidth]{comparison_dynare_vs_chebyshev_part1.png}
    \caption{Сравнение локального  и глобального решений на шок \(\mathit{tfp}\) с \(\sigma_{\text{tfp}} = 0{,}2\): реальные величины и труд --- выпуск \(Y\), капитал \(K\), труд \(L\), потребление \(C\), инвестиции \(I\), зарплата \(W\).}
    \label{fig:local_global_dsge}
\end{figure}

\begin{figure}[htbp]
    \centering
    \includegraphics[width=\textwidth]{comparison_dynare_vs_chebyshev_part2.png}
    \caption{Сравнение локального и глобального решений: цены и факторы --- процентная ставка \(R\), прибыль \(\mathit{Pr}\), предельные издержки \(\mathit{MC}\) и совокупная факторная производительность \(\mathit{tfp}\). Для \(\mathit{Pr}\) показано отклонение в уровнях (стационарное значение равно нулю при функции Кобба-Дугласа), для остальных --- проценты от стационарного состояния.}
    \label{fig:local_global_dsge_2}
\end{figure}

\FloatBarrier

Локальное решение, получаемое через линеаризацию вокруг стационарного состояния, хорошо работает в малой окрестности равновесия и отличается вычислительной простотой. Глобальное решение учитывает нелинейность модели на всем диапазоне состояний и более корректно описывает динамику при больших отклонениях от стационарной точки.

Практический вывод состоит в следующем: для быстрого анализа малых шоков локальные методы удобны, но для задач, связанных с существенными отклонениями и исследованием полной переходной динамики, предпочтителен глобальный подход, который вычислительно дорогой и требует больше времени. 

\paragraph{Монетарный шок: сравнение локального и глобального решений.}

Сравнение проводится по импульсным функциям отклика всех 19 эндогенных переменных на единичный монетарный шок. Для удобства восприятия графики разбиты на два графика - реальные величины и труд (рис.~\ref{fig:mon_compare_real}) и цены, прибыль и факторы (рис.~\ref{fig:mon_compare_prices}).

\begin{figure}[htbp]
    \centering
    \includegraphics[width=\textwidth]{compare_orders_mon_part1.png}
    \caption{Монетарный шок, сравнение локального и глобального решений: реальные величины --- выпуск \(Y\), капитал \(K\), труд \(L\), потребление \(C\), инвестиции \(I\), зарплаты \(W\), \(W^{real}\), аренда капитала \(R\).}
    \label{fig:mon_compare_real}
\end{figure}

\begin{figure}[p]
    \centering
    \includegraphics[width=\textwidth, height=0.92\textheight, keepaspectratio]{compare_orders_mon_part2.png}
    \caption{Монетарный шок, сравнение локального и глобального решений: реальная аренда \(R^{real}\), номинальная ставка \(R^b\), реальная \(R^{b,real}\), предельные издержки \(MC\), \(MC^{real}\), прибыль \(\mathit{Pr}\), \(\mathit{Pr}^{real}\), уровень цен \(P\), инфляция \(\pi\), TFP, монетарный мультипликатор \(\mathit{mon}\).}
    \label{fig:mon_compare_prices}
\end{figure}

\FloatBarrier

Анализ результатов показывает, что для большинства <<номинальных>> переменных (\(W\), \(W^{real}\), \(R\), \(R^{real}\), \(MC\), \(MC^{real}\), \(\mathit{Pr}\), \(\mathit{Pr}^{real}\)) линейное и нелинейное приближения дают существенно различающиеся импульсные отклики. В частности, линейное решение предсказывает крайне резкую начальную реакцию (\(W\) опускается до \(-35\%\), \(R\) --- до \(-80\%\), прибыль вырастает до \(+200\%\)), тогда как нелинейное приближение демонстрирует значительно более умеренные отклики того же знака. Это связано с тем, что параметр \(\mu = 30\) обеспечивает сильную нелинейность в уравнении Ротемберга, которую линеаризация 1-го порядка адекватно не передаёт.

Для <<реальных>> величин (\(Y\), \(K\), \(L\), \(C\), \(I\)) обе схемы дают качественно близкие отклики --- сокращение выпуска и труда, плавное восстановление капитала и потребления, --- но и здесь нелинейное приближение даёт количественно меньшие амплитуды.

Полученные результаты подтверждают общий вывод работы: при анализе реакции экономики на крупные монетарные шоки следует использовать решения более высокого порядка, поскольку линейное приближение систематически переоценивает амплитуду откликов <<номинальных>> переменных.

    \chapter{Выводы}

Выполненная работа была посвящена численному решению модели реального делового цикла для закрытой экономики. В рамках исследования была рассмотрена RBC-модель с репрезентативным домохозяйством и фирмой, выписаны основные условия равновесия, уравнение Эйлера, бюджетное ограничение и соотношения, определяющие стационарное состояние экономики.

Основное внимание в работе уделялось глобальному численному подходу к решению динамической системы. В отличие от локальной линеаризации, которая описывает поведение модели только в малой окрестности стационарного состояния, глобальное решение позволяет анализировать траектории на более широком множестве начальных условий. Это особенно важно для макроэкономических моделей, поскольку при существенных отклонениях от равновесия нелинейные свойства системы начинают играть заметную роль.

В практической части был реализован метод стрельбы, с помощью которого подбиралось начальное потребление, обеспечивающее сходимость траектории к стационарной точке. Для решения неявного уравнения Эйлера на каждом шаге использовался метод Ньютона. Такая схема позволила построить устойчивые переходные траектории капитала, потребления и труда при различных начальных значениях капитала и параметров модели.

Построенные фазовые и трехмерные траектории показали, что экономика сходится к стационарному состоянию как при недостаточном, так и при избыточном начальном уровне капитала. Это подтверждает корректность реализованного численного алгоритма и согласуется с экономической интерпретацией седлово-устойчивой динамики. При этом глобальный подход дает более полное представление о форме переходной траектории, чем локальное приближение.

Сравнение локального и глобального решений показало, что локальные методы удобны для анализа малых отклонений от стационарного состояния, но могут терять точность при больших шоках или далеких начальных условиях. Глобальное решение, напротив, позволяет учитывать особенности модели и ее нелинейности без предварительного сведения системы к линейному виду. Поэтому такой подход является более содержательным при исследовании полной переходной динамики и поведения экономики вне малой окрестности равновесия.

Отдельно следует отметить, что учет нелинейностей открывает возможность анализа эффектов, которые не всегда корректно отражаются в локальной линеаризованной постановке. В частности, в стохастических версиях RBC/DSGE-моделей глобальное решение позволяет выявить эффект сбережений из предосторожности. Его экономический смысл состоит в том, что при наличии неопределенности агенты могут изменять решения о потреблении и накоплении капитала не только из-за среднего ожидаемого шока, но и из-за риска будущих неблагоприятных состояний.

Таким образом, глобальные численные методы дают более гибкий инструмент для анализа макроэкономической динамики. Они позволяют исследовать не только поведение модели около стационарной точки, но и реакцию экономики на существенные отклонения, нелинейные эффекты и потенциальное влияние неопределенности. Полученные результаты могут служить основой для дальнейшего расширения модели, в том числе для учета стохастических шоков, монетарной политики и более детального анализа импульсных откликов.

    \include{mainfiles/chapter5}
    \include{mainfiles/conclusion}

    %Печатаем библиографию
    \clearpage
    \bibliography{diploma}

\end{document}
